\uuid{pQMq}
\titre{Primitives usuelles}
\auteur{hanani ; mzari}
\organisation{exo7}
\licence{CC BY-NC-SA 4.0}
\datecreate{2020-01-20}
\niveau{L1}
\module{Analyse}
\chapitre{Calcul d'intégrales}
\sousChapitre{Primitives diverses}
\difficulte{1}
\theme{Primitives, primitives usuelles}
\exo7qcmid{lille-2.primitives.004}

\contenu{
\begin{question}
\qclassification{127.11}
\qtags{motcle=primitives usuelles}
Parmi les affirmations suivantes, cocher celles qui sont vraies :
\begin{answers}
\good{$\displaystyle F(x)=(x+1)^2+\cos (2x)$ est une primitive de $\displaystyle f(x)=2x+2-2\sin (2x)$ sur $\Rr$.}
\bad{$\displaystyle F(x)=x^2+2x-\cos (2x)$ est une primitive de $\displaystyle f(x)=2x+2-2\sin (2x)$ sur $\Rr$.}
\bad{$\displaystyle F(x)=2x^2\cos (x^2)$ est une primitive de $\displaystyle f(x)=4x\sin (2x)$ sur $\Rr$.}
\good{$\displaystyle F(x)=-2x\cos (2x)+\sin (2x)$ est une primitive de $\displaystyle f(x)=4x\sin (2x)$ sur $\Rr$.}
\end{answers}
\begin{explanations}
On a : $\displaystyle \Big[(x+1)^2+\cos (2x)\Big]'=2x+2-2\sin (2x)$. Donc $\displaystyle (x+1)^2+\cos (2x)$ est une primitive de $\displaystyle 2x+2-2\sin (2x)$ sur $\Rr$. De même, $\displaystyle -2x\cos (2x)+\sin (2x)$ est une primitive de $\displaystyle 4x\sin (2x)$ sur $\Rr$.
\end{explanations}
\end{question}
}
